\uuid{Yc1j}
\titre{Équation différentielle du second ordre et intégrale double}
\niveau{L2}
\module{Analyse}
\chapitre{Equation différentielle}
\sousChapitre{Résolution d'équation différentielle du deuxième ordre}
\theme{Intégrale double, coordonnées polaires}
\auteur{}
\datecreate{2026-06-04}
\organisation{}
\difficulte{3}

\contenu{
\texte{
  Cet exercice comporte deux parties indépendantes.
}

\begin{enumerate}

%%%% PARTIE A %%%%

\item
  \texte{
    \textbf{Partie A — Équation différentielle du second ordre.}\\[0.3em]
    On considère l'équation différentielle :
    \[
      (E') \;:\quad y''(x) - y'(x) - 2\,y(x) = 2x - 1.
    \]
  }
  \question{Résoudre l'équation homogène associée $(H') : y'' - y' - 2y = 0$.}
  \reponse{
    L'équation caractéristique associée est $r^2 - r - 2 = 0$,
    soit $(r-2)(r+1) = 0$, qui admet deux racines réelles distinctes
    $r_1 = 2$ et $r_2 = -1$.
    La solution générale de $(H')$ est donc
    \[
      y_H(x) = c_1\,e^{2x} + c_2\,e^{-x}, \quad c_1, c_2 \in \mathbb{R}.
    \]
  }

\item
  \question{Déterminer une solution particulière de $(E')$
    en cherchant une solution de la forme $y_P(x) = ax + b$.}
  \reponse{
    On pose $y_P = ax + b$, d'où $y_P' = a$ et $y_P'' = 0$.
    En substituant dans $(E')$ :
    \[
      0 - a - 2(ax+b) = 2x-1
      \implies -2ax + (-a-2b) = 2x-1.
    \]
    Par identification :
    \[
      \begin{cases} -2a = 2 \\ -a - 2b = -1 \end{cases}
      \implies
      \begin{cases} a = -1 \\ b = 1. \end{cases}
    \]
    Donc $y_P(x) = -x + 1$.
    \textit{Vérification :} $0 - (-1) - 2(-x+1) = 1 + 2x - 2 = 2x-1$. \checkmark
  }

\item
  \question{En déduire la solution générale de $(E')$ sur $\mathbb{R}$.}
  \reponse{
    La solution générale de $(E')$ est
    \[
      y_G(x) = y_H(x) + y_P(x) = c_1\,e^{2x} + c_2\,e^{-x} - x + 1,
      \quad c_1, c_2 \in \mathbb{R}.
    \]
  }

%%%% PARTIE B %%%%

\item
  \texte{
    \textbf{Partie B — Intégrale double.}\\[0.3em]
    On considère le domaine :
    \[
      D = \left\{(x,y)\in\mathbb{R}^2 \;\middle|\;
          y \geq 0,\quad y \leq x,\quad x^2+y^2 \geq 1,\quad x^2+y^2 \leq 4\right\}.
    \]
  }
  \question{Représenter la région $D$.}
  \reponse{
    Le domaine $D$ est délimité par :
    \begin{itemize}
      \item la droite horizontale d'équation $y = 0$ (axe des abscisses) ;
      \item la droite d'équation $y = x$ ;
      \item le cercle d'équation $x^2+y^2 = 1$, de centre $O$ et de rayon $1$ ;
      \item le cercle d'équation $x^2+y^2 = 4$, de centre $O$ et de rayon $2$.
    \end{itemize}
    $D$ est donc la couronne circulaire $1 \leq r \leq 2$
    intersectée avec le secteur angulaire $0 \leq \theta \leq \pi/4$
    (premier octant, en dessous de $y = x$).
  }

\item
  \question{Calculer $\displaystyle I = \iint_D (x^2+y^2)\;dx\,dy$.
    On pourra utiliser un changement de variables en coordonnées polaires.}
  \indication{Déterminer les bornes de $r$ et $\theta$ correspondant à $D$,
    puis séparer les deux intégrales.}
  \reponse{
    On pose $x = r\cos\theta$, $y = r\sin\theta$, et $dx\,dy = r\,dr\,d\theta$.

    Les contraintes sur $D$ se traduisent en coordonnées polaires :
    \begin{itemize}
      \item $y \geq 0$ et $y \leq x$ donnent $0 \leq \theta \leq \dfrac{\pi}{4}$ ;
      \item $1 \leq x^2+y^2 \leq 4$ donne $1 \leq r \leq 2$.
    \end{itemize}
    Ainsi $D^* = \left\{(r,\theta) \;\middle|\; 0 \leq \theta \leq \dfrac{\pi}{4},\; 1 \leq r \leq 2\right\}$.

    En coordonnées polaires, $x^2 + y^2 = r^2$, donc :
    \begin{align*}
      I &= \iint_{D^*} r^2 \cdot r\,dr\,d\theta
         = \int_0^{\pi/4} d\theta \int_1^2 r^3\,dr \\
        &= \frac{\pi}{4} \cdot \left[\frac{r^4}{4}\right]_1^2
         = \frac{\pi}{4} \cdot \frac{16-1}{4}
         = \frac{15\pi}{16}.
    \end{align*}
  }

\end{enumerate}
}
